\section{Contributing clocks and the entropy of the profile}
\label{fp:sec:arithmetic}

We apply the neighboring-moment estimates to one fixed accepting source
from Sections~\ref{lower:chapter}--\ref{lower:frozen}. All its original
physical groups, fine types, count windows, allocated walls, costs and
adaptive targets are retained. The source measure below includes the
regularity and bank acceptance on unselected prefix blocks as well as
the selected binary comparisons.

\subsection{One source for a compact interval of moment orders}

Fix a dimension-dependent compact interval
$J=[p_-,p_+]\Subset(0,1)$ which contains the whole range of owner
products in \eqref{profile:compact} in its interior. The margin from
these products to the endpoints of $J$ is fixed and positive. The
physical saddle and every original mark law are held fixed when a
test order $p\in J$ is varied.

For a physical vector $L$ in the sufficiently large positive-length
region, let $\mathcal I$ be the selected native binary comparisons
of Section~\ref{lower:data}. There is at most one such comparison
in an original physical block. If $i\in\mathcal I$, write $e_i$
for its physical canonical edge and put
\begin{equation}
 \eta_i=\sigma_{\operatorname{owner}(i)},\qquad
 \rho_i=\rho_{e_i},\qquad
 \tau_i=-2\log\min\{1,T_{e_i}\},
 \label{fp:eq:physical-duration}
\end{equation}
where $T_e,\rho_e$ are exactly those of
\eqref{profile:edge-factor}. In particular
$P_{e_i}=\exp\{-\tau_i\chi(\eta_i,\rho_i)\}$.
The durations and correlations in this section are deterministic
functions of the physical vector. Every sum involving them is over
selected comparisons unless another index set is displayed.
Unselected blocks have $D_i=1$ and contribute no duration or Gaussian
parameter; they still retain their actual local acceptance.

\begin{lemma}[Uniform neighboring moments of a fixed cell source]
\label{fp:lem:fixed-source-window}
The complete-cell source and its costs can be chosen once, uniformly
for $p\in J$, so that at every supported complete prefix type
$\omega$ and admitted original count vector,
\begin{equation}
 \begin{split}
 &0\le q_i(\omega,g_i)\le1,\qquad
 0<D_i(g_i)\le1,\qquad q_i\le C_kD_i,\\
 &\EE_{g_i}q_i(\omega,g_i)^p
   \asymp_k\EE_{g_i}D_i(g_i)^p
   \asymp_k\exp\{-\tau_i\chi(p,\rho_i)\},\qquad p\in J.
 \end{split}
 \label{fp:eq:actual-window}
\end{equation}
Here $g_i$ has the original exact-total cell law
\eqref{ld:exact-cell-law}. The same $q_i,D_i$ occur for every $p$.
Each $D_i$ reads only the uncolored cell occupancies, including their
total, and is independent of the fine type after that total is fixed.
\end{lemma}
\begin{proof}
Choose the compact scalar ranges with the test order ranging over
all of $J$. Theorem~\ref{bridge:all-reserves} and
Corollary~\ref{bridge:contracted} then supply common wall, reserve
and binary regularity constants. The refinement in
Lemma~\ref{ld:named-REG} reads both true mark lists and has a uniform
conditional success probability for each retained binary word and
full type. Choose its constants once for these same walls.

In \eqref{ld:cost-envelope}, choose a single floor exponent satisfying
$D_0p_->2\sup_{p\in J,\rho}\chi(p,\rho)+1$ over the compact
correlation range. Fix the literal exit widths and use this floor to
bound the whole logarithmic target and recoloring range. Then choose
the coefficient tail and the all-transfer envelopes
\eqref{ld:all-e-parameters} uniformly for $p\in J$, with their
geometric depth coefficient fixed below the available history budget.
Finally choose the spare strong-wall and named regularity margins,
macroscopic-bank support, complete-cell allowances, and final count
and root thresholds, in the order of \eqref{ld:targets} and the end
of Section~\ref{ld:cells}. None of these later choices changes $D_0$
or requires evaluating a target before its envelope has been fixed.

The lower bridge moment, the upper envelope moment
\eqref{ld:cost-moment}, and the relative macroscopic-bank retention
\eqref{ld:local-bank-retention} are uniform in this compact order
range. The direct cell construction \eqref{ld:cell-local} preserves
them directly for the same functions of the cell word, with the
private acceptance evaluated before taking its power. The pointwise inequality
$q_i\le C_kD_i$ supplies the other sides of the two moment
comparisons.

Initially these comparisons use the exact count, compatible local
pair and centered correlation of \eqref{lower:conditional-factor}.
The central-count estimate \eqref{ld:count-comparison} and the
physical-correlation comparison following
\eqref{lower:natural-correlation} are uniform in $p\in J$ by the
compact Lipschitz bound for $\chi$. More explicitly, on a rare
branch the additional correlation error is
$O_k(T_{e_i}+L_i^{-1})$ and
$T_{e_i}\ge c_k/(1+L_i)$; its product with $|\log T_{e_i}|$ is
bounded. Replacing endpoint scales by fixed comparable scales changes
the effective duration by $O_k(1)$. Threshold crossings and
saturated comparisons also have bounded effective duration.
These observations yield \eqref{fp:eq:actual-window} with the
physical parameters \eqref{fp:eq:physical-duration}.

On the actual arithmetic grid, keep the physical saddle frozen as
in Section~\ref{lower:frozen}. Its bounded grid defects and exact
local increments give the same endpoint and
exponent-times-log comparisons uniformly over $J$. The preceding
argument therefore applies directly to the grid functions. The
resulting neighboring-moment inequalities determine the physical
centers below. Bounded supported experiments retain their
prescribed finite support and fixed positive mass throughout.
\end{proof}

\subsection{The whole prefix source and the unselected blocks}

Let $\mathcal P$ be the set of original blocks before the unread
ordinary suffix. For each admitted prefix total vector $\mathbf n$,
let $\mu_{\mathbf n}$ be the one count-only measure on complete
prefix types in \eqref{ld:whole-count-mass}. Its mass is uniformly
comparable to $\prod_{j<m}w_j$. Conditional on a complete type,
all original cell clocks are independent. At a selected block retain
$q_i$ from Lemma~\ref{fp:lem:fixed-source-window}. At
$i\in\mathcal P\setminus\mathcal I$, denote its actual private
local regularity and bank acceptance by $q_i^{\rm reg}(\omega,g_i)$.
The unselected-block case of \eqref{ld:local-bank-retention},
together with the named refinement, gives
\begin{equation}
 0\le q_i^{\rm reg}\le1,\qquad
 c_k\le\EE_{g_i}q_i^{\rm reg}(\omega,g_i)\le1
 \label{fp:eq:unselected-mass}
\end{equation}
uniformly in the whole supported type. Equivalently, the lower
fractional moment already established with $D_i=\Psi_i=1$ implies
this lower first moment by Jensen.

To include the remaining central uncolored totals explicitly, let
$\nu(d\mathbf n)$ be their original independent natural Poisson
law restricted to their prescribed central windows. Put
\begin{equation}
 \mathcal W=\int\mu_{\mathbf n}(1)\,\nu(d\mathbf n),\qquad
 d\PP_{\rm ref}
 =\mathcal W^{-1}\nu(d\mathbf n)\mu_{\mathbf n}(d\omega)
       \prod_{i\in\mathcal P}d\PP_{i,n_i}(g_i),
 \label{fp:eq:source-reference}
\end{equation}
where $\PP_{i,n_i}$ is \eqref{ld:exact-cell-law}. The same
definitions can be used conditional on any one admitted total
vector, by taking $\nu$ concentrated there. This reference law
retains the coupled type coordinates under the count-only acceptance.

Set
\begin{equation}
 \begin{split}
 U(\omega,g)&=\prod_{i\in\mathcal P\setminus\mathcal I}
                      q_i^{\rm reg}(\omega,g_i),\\
 R_i&=q_i/D_i,\qquad
 W_{\rm pre}(\omega,g)
       =U(\omega,g)\prod_{i\in\mathcal I}R_iD_i^{\eta_i},\\
 Z_{\rm pre}&=\EE_{\rm ref}W_{\rm pre},\qquad
 d\mathfrak M=Z_{\rm pre}^{-1}W_{\rm pre}\,d\PP_{\rm ref}.
 \end{split}
 \label{fp:eq:source-weight}
\end{equation}
The deterministic factor $\prod_{j<m}w_j^{\sigma_j-1}$ from
the owner costs cancels on normalization. Thus $\mathfrak M$ is
the complete prefix mixed source, including its unselected blocks.

\begin{lemma}[Integrating the unselected acceptance]
\label{fp:lem:unselected-integration}
At a fixed $(\mathbf n,\omega)$, integrating the unselected
clocks in \eqref{fp:eq:source-weight} replaces their factor by
\begin{equation}
 H(\mathbf n,\omega)
   =\prod_{i\in\mathcal P\setminus\mathcal I}
                     \EE_{g_i}q_i^{\rm reg}(\omega,g_i),
 \qquad c_k\le H\le1.
 \label{fp:eq:unselected-type-weight}
\end{equation}
This is a bounded reweighting of the same coupled count-only type
measure. In particular, uniformly on compact subsets of
$s_i<1$, $p_i=\eta_i-s_i-t_i\in\operatorname{int}J$,
\begin{align}
 Z_{\rm pre}&\asymp_k
       \exp\{-\sum_{i\in\mathcal I}\tau_i\chi(\eta_i,\rho_i)\},
       \label{fp:eq:complete-source-Z}\\
 \log\EE_{\mathfrak M}
       e^{\sum_{i\in\mathcal I}(s_iS_i+t_iT_i)}
   &=\sum_{i\in\mathcal I}\tau_i
       [\chi(\eta_i,\rho_i)-\chi(p_i,\rho_i)]+O_k(1),
       \label{fp:eq:complete-source-spectrum}
\end{align}
where $S_i=-\log q_i$, $T_i=-\log D_i$ on the mixed support.
\end{lemma}
\begin{proof}
At a fixed complete type, the unselected factors use only their
own original clocks. Their exact integral is the product in
\eqref{fp:eq:unselected-type-weight}; its bounds follow from
\eqref{fp:eq:unselected-mass} and the fixed number of blocks.
The selected clocks remain independent conditional on that type.
For arbitrary fixed $a_i>0$ and $p_i\in J$, Lemma~\ref{fp:lem:mixed-power}
and \eqref{fp:eq:actual-window} give
\[
 \EE_{g_i}R_i^{a_i}D_i^{p_i}
      \asymp_k e^{-\tau_i\chi(p_i,\rho_i)}.
\]
Multiply these integrals at the fixed type, retain $H$ inside the
whole type sum, and use its uniform upper and lower bounds. Tonelli
then integrates that one coupled measure and the central totals.
Taking $a_i=1,p_i=\eta_i$ gives the normalization. For the moment
numerator take $a_i=1-s_i$ and use the exact identity
$R_iD_i^{\eta_i}e^{s_iS_i+t_iT_i}=R_i^{1-s_i}D_i^{p_i}$.
Division by the normalization proves the spectrum with $O_k(1)$
error.
\end{proof}

\subsection{Clock shells, entropy, and projection}

Define physical centers and an uncolored prefix event by
\begin{equation}
 m_i=\tau_i\partial_p\chi(\eta_i,\rho_i),\qquad
 \mathcal G_A=\bigcap_{i\in\mathcal I}
       \{|T_i-m_i|\le A\sqrt{1+\tau_i}\}.
 \label{fp:eq:source-shell}
\end{equation}

\begin{proposition}[Contributing clock shells]
\label{fp:prop:source-shell}
Uniformly over the physical large-length region and the admitted
central-count support,
\begin{equation}
 \mathfrak M(\mathcal G_A^c)\le C_ke^{-c_kA},\qquad A\ge1.
 \label{fp:eq:shell-tail}
\end{equation}
Both $S_i$ and $T_i$ satisfy
\begin{equation}
 \mathfrak M\{|S_i-m_i|>u\},\quad
 \mathfrak M\{|T_i-m_i|>u\}
 \le C_k\exp\{-c_k\min(u^2/(1+\tau_i),u)\}.
 \label{fp:eq:complete-source-concentration}
\end{equation}
Moreover
\begin{equation}
 \mathfrak M\{\log(D_i/q_i)>v\}\le C_ke^{-v},\qquad
 \log(D_i/q_i)\ge-\log C_k,
 \label{fp:eq:complete-source-ratio}
\end{equation}
and
\begin{equation}
 \KL(\mathfrak M\Vert\PP_{\rm ref})
  =\sum_{i\in\mathcal I}\tau_i
       [\chi(\eta_i,\rho_i)-\eta_i\partial_p\chi(\eta_i,\rho_i)]
       +O_k\!\left(\sum_{i\in\mathcal I}\sqrt{1+\tau_i}+1\right).
 \label{fp:eq:source-entropy}
\end{equation}
The same entropy formula holds for projection to the complete
uncolored prefix occupancies, with the corresponding reference
marginal. The expected conditional entropy of the discarded type
information is $O_k(1)$. The event $\mathcal G_A$ reads neither a
fine name nor an unread ordinary-suffix clock.
\end{proposition}
\begin{proof}
The local $C^{1,1}$ estimate of Theorem~\ref{fp:thm:regularity},
applied to \eqref{fp:eq:complete-source-spectrum}, gives the
two-sided Bernstein bound by Chernoff exactly as in
Proposition~\ref{fp:prop:history-target}. Here the order interval
has a fixed margin around every $\eta_i$. For $A\ge1$,
$\min(A^2,A\sqrt{1+\tau_i})\ge A$; a union bound gives the shell
estimate. At a fixed type, the proof of the ratio tail uses
$R_i\mathbf1_{\{R_i<e^{-v}\}}\le e^{-v}$. The other selected
integrals and the factor $H$ have uniform comparisons, so the same
argument proves \eqref{fp:eq:complete-source-ratio} for the
complete source. In particular each $\EE_{\mathfrak M}|\log R_i|$
is bounded.

Conditional on the type and selected clocks, the unselected clocks
have density $U/H$ relative to their original product law. Since
$0\le U\le1$ and $H\ge c_k$,
\[
 \EE_{\mathfrak M}[-\log U\mid\mathbf n,\omega,g_{\mathcal I}]
     =H^{-1}\EE[U(-\log U)\mid\mathbf n,\omega]
     \le (e c_k)^{-1}.
\]
Thus the unselected source restrictions have bounded entropy cost.
Integrating \eqref{fp:eq:complete-source-concentration} gives
$\EE_{\mathfrak M}T_i=m_i+O_k(\sqrt{1+\tau_i})$.
Expansion of the exact density
\[
 \log\frac{d\mathfrak M}{d\PP_{\rm ref}}
  =\log U+\sum_{i\in\mathcal I}\log R_i
       -\sum_{i\in\mathcal I}\eta_iT_i-\log Z_{\rm pre}
\]
proves \eqref{fp:eq:source-entropy}, including the case
$\mathcal I=\varnothing$.

For clarity, define the two Gibbs laws relative to this same full
reference measure by densities proportional to
$\prod_{i\in\mathcal I}D_i^{\eta_i}$ and
$\prod_{i\in\mathcal I}q_i^{\eta_i}$, respectively. Their
normalizers and $Z_{\rm pre}$ are comparable. The respective
density ratios of $\mathfrak M$ are bounded by constants times
$U\prod_iR_i$ and $U\prod_iR_i^{1-\eta_i}$, hence are uniformly
bounded. Both forward Gibbs relative entropies, and the corresponding
finite-time Gibbs variational excesses, are consequently $O_k(1)$.
For the $D$-Gibbs law its density is measurable with respect to the
complete uncolored occupancies (including the totals). Its conditional
type law given those occupancies agrees exactly with that of
$\PP_{\rm ref}$. The relative-entropy chain rule, as in
Corollary~\ref{fp:cor:projection}, bounds the entropy of the
discarded information by $O_k(1)$ and proves the marginal formula.
All physical centers remain deterministic when the original central
totals are integrated in \eqref{fp:eq:source-reference}.
\end{proof}

\subsection{Restriction and transfer to original prime tuples}

\begin{theorem}[Full-order arithmetic lower family]
\label{fp:thm:arithmetic-family}
For each fixed $k$, the source constants and complete-cell conventions
can be chosen as above. There are fixed $L_*(k)$ and $A(k)$ such
that for every physical vector $L_i\ge L_*(k)$ there is a family
$\mathcal A_A(L)$ of the original squarefree tuples in
\eqref{paper:K} satisfying
\begin{equation}
 e^{-\sum_i h_iL_i}
 \sum_{\mathbf a\in\mathcal A_A(L)}
       \frac{L^{(k+1)}(\mathbf a)}{a_1\cdots a_k}
       \asymp_k\Psi_k(L),\qquad h_i=k-i+2.
 \label{fp:eq:arithmetic-family}
\end{equation}
Every tuple in this family has complete-cell occupancy in
$\mathcal G_A$ and retains the original uncolored count/root support
and size restrictions. Thus the family contributes a fixed positive
dimension-dependent fraction of the entire harmonic volume sum.
\end{theorem}
\begin{proof}
The joint ordinary-suffix kernel \eqref{ld:cell-root-kernel}
holds pointwise on each retained complete prefix. Multiply it by
the original owner cost factors and
$\mathbf1_{\mathcal G_A^c}$. The complete prefix measure in
\eqref{fp:eq:source-weight}, including $U$, is exactly the one
which results when these factors are integrated against the private
prefix mass. Equations~\eqref{fp:eq:shell-tail} and
\eqref{ld:cell-root-kernel} therefore bound the removed mixed
source mass by $C_ke^{-c_kA}$ times the original mixed source mass.

Keep the costs, floors, targets, strong curves, count bands and root
rule unchanged, and restrict the accepted source by
$\mathbf1_{\mathcal G_A}$. The recoloring kernel preserves every
uncolored cell occupancy, so its transported density is exactly
\begin{equation}
 F_A(g,\omega)=\mathbf1_{\mathcal G_A}(g)F_y^{I,\rm cell}(g,\omega).
 \label{fp:eq:restricted-density}
\end{equation}
The pointwise mass identity \eqref{ld:cell-density} includes its
original root indicator. In the exact nominal identity
\eqref{ld:exact-first-mass}, the remaining factors involving
$E_y$ and the root error are bounded above and below. Consequently
the removed nominal first mass is at most $C_ke^{-c_kA}$ of the
whole first mass. Choose the fixed $A(k)$ large enough to retain
a fixed positive fraction. The nominal mass remains comparable to
$\Psi_k(L)$ by \eqref{lower:nominal-mass} and the frozen-grid
comparison of Section~\ref{lower:frozen}.

For any retained anchor and any original history,
$M_h(F_A;\omega)\le M_h(F_y^{I,\rm cell};\omega)$. The original
unnormalized collision row \eqref{lower:unnormalized-row} thus
retains its constant under the restriction. For every retained
occupancy the pointwise squarefree supply
\eqref{lower:actual-supply} also remains valid. Apply the original
prime-cell kernel to $F_A$ on that occupancy. With
$Z_{F_A}(g)=\sum_\omega F_A(g,\omega)$, it gives
\[
 \int f_A\ge c_kW_gZ_{F_A}(g),\qquad
 \int f_A^{1+\delta}\le C_kW_gZ_{F_A}(g).
\]
The same H\"older inequality as \eqref{lower:holder} yields
$\operatorname{meas}\{f_A>0\}\ge c_kW_gZ_{F_A}(g)$.

Here is an explicit tuple definition. For every occupancy $g$ let
$\mathcal A(g)$ be its original squarefree tuple set from
\eqref{lower:actual-supply}, with the source size restrictions, and set
\begin{equation}
 \mathcal A_A(L)=
    \bigcup_{g:Z_{F_A}(g)>0}\mathcal A(g).
 \label{fp:eq:tuple-family-definition}
\end{equation}
These sets have disjoint occupancies. The positive set of $f_A$ is
contained in the original divisor union for each such tuple.
Summing the preceding bound over $g$ and retaining the nominal
normalization proves the lower bound in
\eqref{fp:eq:arithmetic-family}. Its upper bound follows by
inclusion in the entire sum \eqref{paper:K} and
Theorem~\ref{positive:upper}. Bounded grid-length losses change the
normalizing exponent only by a dimension-dependent constant.

All these objects are constructed on the actual complete prime-cell
grid with the physical products and owner structure frozen. The
grid comparison \eqref{lower:frozen-defects} uses the exact budget
of Section~\ref{lower:frozen}.
Since \eqref{fp:eq:restricted-density} only removes accepted words,
the original count, root, occupancy and size support is preserved.
\end{proof}

Each $\mathbf a\in\mathcal A_A(L)$ determines its occupancies
$g_{i,c}(\mathbf a)$ by counting its prime factors in the fixed
complete cells. We define $D_i(\mathbf a)=D_i(g_i(\mathbf a))$.
The shell and targets are therefore actual functions of the original
tuple, evaluated from its uncolored occupancies. The condition
$Z_{F_A}(g)>0$ in \eqref{fp:eq:tuple-family-definition} expresses
the existence of an accepted fine word; $q_i(\omega,g_i)$ still
depends on that word's fine type.

\subsection{The original adaptive targets and their physical centers}

Retain \eqref{lower:owner-cost}, \eqref{lower:targets} and
\eqref{lower:triangular}, with $w_m=1$ for the terminal owner:
\[
 y_j=-\log w_j+
       \sum_{\substack{i\in\mathcal I\\\operatorname{owner}(i)=j}}T_i,
 \qquad x=\sum_j\frac{\sigma_j}{\sigma_1}y_j.
\]
Define deterministic physical centers by
\begin{align}
 \bar y_j&=-\log w_j+
       \sum_{\substack{i\in\mathcal I\\\operatorname{owner}(i)=j}}
           \tau_i\partial_p\chi(\sigma_j,\rho_i),
           \label{fp:eq:ybar}\\
 \bar x&=\sum_j\frac{\sigma_j}{\sigma_1}\bar y_j.
           \label{fp:eq:xbar}
\end{align}

\begin{corollary}[Target and root localization on original tuples]
\label{fp:cor:targets}
Every tuple in \eqref{fp:eq:tuple-family-definition} satisfies
\begin{align}
 |y_j-\bar y_j|&\le C_kA
       \sqrt{1+\sum_{\substack{i\in\mathcal I\\
                                      \operatorname{owner}(i)=j}}\tau_i},
       \label{fp:eq:y-error}\\
 |x-\bar x|&\le C_kA\sqrt{1+\sum_{i\in\mathcal I}\tau_i}.
       \label{fp:eq:x-error}
\end{align}
Let $L'_i$ be the actual complete-cell lengths,
$V'_j=\sum_{i\le j}L'_i$, and retain the physical $s_i,G_i$.
The exact uncolored grid root on a tuple is
\begin{equation}
 X_1^{\rm grid}(\mathbf a)=
    \frac{\sum_iG_i\,\omega(a_i)-\sum_js_jV'_j}{\sigma_1},
 \qquad
 X_1^{\rm grid}=x+\epsilon,\qquad -C_k\le\epsilon\le0.
 \label{fp:eq:tuple-root}
\end{equation}
Here $\omega(a_i)$ denotes the number of prime factors of the
squarefree integer $a_i$. If $e_j$ is the original integer recoloring
recursion, then it differs by
$O_k(A\sqrt{1+\sum_i\tau_i}+1)$ from the deterministic real recursion
\[
 \bar e_0=0,\qquad
 \bar e_j=(\bar y_j+A_j\bar e_{j-1})/b_j.
\]
\end{corollary}
\begin{proof}
Sum the shell inequalities within each owner and use Cauchy--Schwarz
and the dimension-bounded number of selected blocks. The ratios
$\sigma_j/\sigma_1$, the reciprocals of $b_j$, and the remaining
recursion coefficients are dimension-bounded. Iterating the
recursion difference, with the error in $[0,1)$ at each ceiling,
proves the stated recoloring estimate.

The first expression in \eqref{fp:eq:tuple-root} is the uncolored
count/root identity \eqref{paper:ep-root-telescope} on the complete
original chain, evaluated algebraically at the frozen physical
coefficients and the actual grid budget.
The actual reservoir count is the floor in \eqref{ld:new-root},
using this grid budget as specified in Section~\ref{lower:frozen}.
Thus its floor error is in $[-r_m w_*,0]$, where
$r_m=\sigma_m/\sigma_1$ and $w_*\le\log(k+1)$ is the reservoir
mark. This proves the second assertion in
\eqref{fp:eq:tuple-root}. Recoloring fixes all uncolored counts,
and the shell restriction fixes the root rule, so the identity
holds on every tuple of the family. Thus the root is the complete-cell
count functional with frozen marks $G_i$ and grid budgets $V_j'$.
\end{proof}

The derivative in \eqref{fp:eq:ybar} is with respect to the test
moment order while the physical vector, mark law and correlation
are fixed. The exact support and recoloring rules continue to use
the adaptive targets $y_j,x$; the centers describe their localization.

For the complete normalized prefix source, every fixed real linear
path target $X_{\rm path}=\sum_{i\in\mathcal I}a_iT_i$ also obeys
\begin{equation}
 \log\EE_{\mathfrak M}e^{hX_{\rm path}}
   =\sum_{i\in\mathcal I}\tau_i
       [\chi(\eta_i,\rho_i)-\chi(\eta_i-ha_i,\rho_i)]+O_k(1)
 \label{fp:eq:source-target-spectrum}
\end{equation}
for small fixed $h$. A deterministic count-window contribution can
be added exactly; the conditional statement likewise holds at fixed
admitted totals. The selected local deviation conclusion of
Theorem~\ref{fp:thm:history-local-ldp} applies when the effective
durations have a common diverging scale.

\begin{remark}[Scope of the arithmetic localization]
\label{fp:rem:arithmetic-scope}
The localization describes a full-order subfamily in the sufficiently
large positive-length region, measured by harmonic divisor volume.
The local deviation principle concerns the normalized prefix source.
For bounded or zero physical lengths, arithmetic stability controls
the total sum; localization of the contributing clocks in that region
would require a separate argument.
\end{remark}

\subsection{Environmental entropy and the structural profile}
\label{fp:sec:interpretation}
The scalar exponent measures a balance between selecting a favorable
clock and asking its private labels to survive. To make this balance
precise, let $\mathsf P_t$ be the OU environmental law and write
$S_t=-\log q_t$. The classical Gibbs variational formula
\cite[Section~1.4]{DupuisEllis1997}, applied to $\eta S_t$, gives
\begin{equation}\label{fp:eq:gibbs-identity}
 b_t(\eta,\rho)=\inf_{\mathsf Q\ll\mathsf P_t}
 \{\KL(\mathsf Q\Vert\mathsf P_t)+\eta\EE_{\mathsf Q}S_t\}.
\end{equation}
Indeed, for $d\nu_{\eta,t}=q_t^\eta d\mathsf P_t/F_{\eta,\rho}(t)$,
subtracting $b_t$ from the expression in braces gives
$\KL(\mathsf Q\Vert\nu_{\eta,t})$. Define
\begin{equation}\label{fp:eq:selection-cost}
 v(\eta,\rho)=\partial_\eta\chi(\eta,\rho),\qquad
 d(\eta,\rho)=\chi(\eta,\rho)-\eta v(\eta,\rho).
\end{equation}
Theorem~\ref{fp:thm:regularity} gives
\[
 \EE_{\nu_{\eta,t}}S_t=t v(\eta,\rho)+O(\sqrt t),\qquad
 \KL(\nu_{\eta,t}\Vert\mathsf P_t)
       =t d(\eta,\rho)+O(\sqrt t).
\]
In particular $d(\eta,\rho)\ge0$. The complete mixed source has the
corresponding sums of these leading quantities by
Proposition~\ref{fp:prop:source-shell}. Its selected local factors
are $qD^{\eta-1}$, and its unselected acceptance has bounded entropy
cost.

The full scalar large-deviation principle in
Theorem~\ref{fp:thm:scalar-ldp} makes the environmental cost explicit.
Its good rate function $I_\rho$ satisfies
\begin{equation}\label{fp:eq:scalar-var}
 \begin{split}
 \chi(\eta,\rho)&=\min_{v\ge0}\{I_\rho(v)+\eta v\},\\
 \argmin_{v\ge0}\{I_\rho(v)+\eta v\}
   &=\{\partial_\eta\chi(\eta,\rho)\}.
 \end{split}
\end{equation}
These identities follow from convex duality for
$\Lambda_\rho(\theta)=-\chi(-\theta,\rho)$ at $\theta=-\eta$;
differentiability makes its subgradient a singleton. Thus
$d(\eta,\rho)=I_\rho(v(\eta,\rho))$.

For a canonical edge $e$ in \eqref{profile:edge-factor}, write
$\sigma_e=\sigma_{\operatorname{owner}(e)}$ for its owner product and set
\begin{equation}\label{fp:eq:edge-clock}
 \tau_e=2[-\log T_e]_+,\qquad
 P_e=e^{-\tau_e\chi(\sigma_e,\rho_e)}.
\end{equation}
Proposition~\ref{fp:prop:poisson} and \eqref{fp:eq:scalar-var} now
express the positive-length profile as two entropy costs:
\begin{align}\label{fp:eq:profile-var}
 -\log\Psi_k(L)
 ={}&\sum_iL_i\mathcal Q(Z_i)-\log\mathcal B(L)
       +\sum_{j\ \mathrm{proper}}
          \sigma_j\log(1+d_j\sqrt{1+T_j})\notag\\
 &+\inf_{(v_e\ge0)}\sum_e\tau_e
             \{I_{\rho_e}(v_e)+\sigma_ev_e\}.
\end{align}
Here $T_j$ denotes owner $j$'s native length as in
\eqref{profile:raw}. For every active edge the minimizing coordinate is
$v_e=\partial_\eta\chi(\sigma_e,\rho_e)$; a coordinate with
$\tau_e=0$ has no effect. The sum separates by the proved product
formula \eqref{profile:raw}: each summand uses the scalar OU rate
function for its edge.
The Poisson cost is measured in the physical lengths $L_i$, whereas
the persistence cost is measured in the effective times $\tau_e$,
which are often logarithmic in those lengths. The constant-factor
bridge estimates remain necessary to recover the arithmetic order.

\subsection{The loss from averaging the clock first}
Let $\Psi_k^{\mathrm{ann}}$ be the expression \eqref{profile:raw}
with each $\chi(\sigma_e,\rho_e)$ replaced by $\sigma_e/2$.
This substitution keeps the saddle, endpoint scores, count windows
and ordinary suffix fixed.
\begin{corollary}\label{fp:cor:anneal-profile}
Uniformly on the positive-length domain,
\begin{equation}\label{fp:eq:anneal-profile}
 \log\frac{\Psi_k^{\mathrm{ann}}(L)}{\Psi_k(L)}
 =\sum_e\tau_e\left[\chi(\sigma_e,\rho_e)-\frac{\sigma_e}{2}\right]
 \asymp_k\sum_e\tau_e.
\end{equation}
Both sides are zero if all $\tau_e$ vanish. Along a parameter family,
the ratio tends to infinity exactly when $\sum_e\tau_e\to\infty$.
\end{corollary}
\begin{proof}
The owner products are bounded away from zero and one. On the
canonical nondegenerate parameter range, the correlations have a
positive lower bound depending only on $k$. Apply
Theorem~\ref{fp:thm:annealing} for the lower estimate and
$\chi(\sigma_e,\rho_e)\le1/2$ for the upper estimate.
\end{proof}

\subsection{The critical six-coordinate family}
In the family \eqref{ex:critical-lengths}, take the fixed middle
length $C$ above the threshold of Theorem~\ref{fp:thm:arithmetic-family}.
The rare edge is $6\to1$, with owner product $\sigma_c$ and correlation
\[
 \rho_c=\frac{(7\log7-2\log2)^2}
              {(7\log7-2\log2)^2+10(\log2)^2}.
\]
Its low-mark probability is $2/7$, so this is exactly
$\rho(6,1;1)$ in \eqref{ex:rho}. Put
\[
 \chi_c=\chi(\sigma_c,\rho_c),\qquad
 v_c=\partial_\eta\chi(\sigma_c,\rho_c),\qquad
 d_c=\chi_c-\sigma_cv_c=I_{\rho_c}(v_c).
\]
The endpoints \eqref{ex:critical-endpoints} give
$\tau_e=\log T+O_k(1)$. Thus \eqref{ex:critical-answer} becomes
\begin{equation}\label{fp:eq:six-example}
 \Phi_6(\log3+L)
 \asymp_6 e^{-J(L)}T^{-1/2}\,T^{-d_c}\,T^{-\sigma_cv_c}.
\end{equation}
The last two factors respectively describe the environmental selection
cost and the conditional path cost at the selected clock. Under the
mixed source, the environmental entropy is
$d_c\log T+O_k(\sqrt{\log T})$. Moreover, the arithmetic family in
Theorem~\ref{fp:thm:arithmetic-family} can be chosen so that
\begin{equation}\label{fp:eq:six-shell}
 -\log D=v_c\log T+O_k(\sqrt{\log T}),\qquad
 X_1^{\mathrm{grid}}=v_c\log T+O_k(\sqrt{\log T}).
\end{equation}
Here $X_1^{\mathrm{grid}}$ is the uncolored count root of the
complete-cell construction, including its bounded rounding error,
as in Corollary~\ref{fp:cor:targets}. The bounded clocks contribute
only bounded target terms. This locates a full-order lower family
of original prime tuples.
Finally the comparison that averages first is too large by a factor
$\asymp T^{\Delta_c}$, where
\[
 \Delta_c=\chi_c-\sigma_c/2
 \ge\frac{\rho_c}{16\pi}\min\{\sigma_c,1-\sigma_c\}>0.
\]
Since $T=\log(\log y_1/\log3)$, this is a power of $\log\log y_1$.

\subsection{How overlap changes the environmental cost}
The original comparison constraints can share a clock. A separate
calculation for deterministic OU windows shows explicitly why the
cost of a shared environment differs from that of independent ones.
Use the conditional probabilities $q_{j,I}$ and joint pressures
$\Xi_A$ of \eqref{fp:eq:joint-pressure}.
Theorem~\ref{fp:thm:overlap} gives
\[
 \EE\prod_j q_{j,I_j}^{p_j}
 \asymp\exp\left[-\int_\R
       \Xi_{\{j:s\in I_j^\circ\}}(\boldsymbol p;\boldsymbol\beta)\,ds\right]
\]
for any fixed number of deterministic intervals, uniformly in their
lengths and positions. If all slopes equal $\sqrt{\rho/(1-\rho)}$,
then the conditional probabilities on the same window coincide, so
\begin{equation}\label{fp:eq:overlap-scalar}
 \Xi_A=\chi\left(\sum_{j\in A}p_j,\rho\right),\qquad
 \EE\prod_jq_{\rho,I_j}^{p_j}
 \asymp\exp\left[-\int_\R\chi(w_{\boldsymbol p}(s),\rho)\,ds\right],
\end{equation}
where $w_{\boldsymbol p}=\sum_jp_j\ind_{I_j}$ and the value
$\chi(0,\rho)=0$ denotes the region with no active interval.
For $I_j=t[a_j,b_j]$, write
$\Xi_\Gamma(\boldsymbol p)=\int\chi(\sum_jp_j\ind_{[a_j,b_j]}(s),\rho)\,ds$.
The finitely many active sets and Theorem~\ref{fp:thm:regularity} give
\begin{equation}\label{fp:eq:overlap-derivative}
 \partial_{p_j}\Xi_\Gamma
 =\int_{a_j}^{b_j}\partial_p\chi(w_{\boldsymbol p}(s),\rho)\,ds,
\end{equation}
where in this formula $w_{\boldsymbol p}=\sum_jp_j\ind_{[a_j,b_j]}$.
The entropy rate of the joint Gibbs environment is consequently
\begin{equation}\label{fp:eq:overlap-entropy}
 \Xi_\Gamma-\boldsymbol p\cdot\nabla\Xi_\Gamma
 =\int\{\chi(w_{\boldsymbol p}(s),\rho)
       -w_{\boldsymbol p}(s)\partial_p\chi(w_{\boldsymbol p}(s),\rho)\}\,ds,
\end{equation}
with zero integrand outside the active union. The same leading rate
holds for mixed envelopes satisfying the joint window
\eqref{fp:eq:joint-window}.

For two intervals of length $t$ with overlap $h$, the formula reads
\begin{equation}\label{fp:eq:overlap-two}
 \EE[q_{I_1}^{p_1}q_{I_2}^{p_2}]
 \asymp e^{-(t-h)[\chi(p_1,\rho)+\chi(p_2,\rho)]
                  -h\chi(p_1+p_2,\rho)}.
\end{equation}
In particular, at $p_1=p_2=1/2$ it equals, up to constants,
$e^{-2(t-h)\chi(1/2,\rho)-h/2}$. These deterministic-window
identities explain the shared-clock effect; the arithmetic profile
uses the conditioned, exact-type construction proved above.
